%%% Discrete Lyapunov Stability of SGD in the Interpolation Regime:
%%% Per-Sample Analysis and Generalization Bounds
%%%
%%% A self-contained discrete-time analysis.
%%% No SDE approximation; all results hold directly for the SGD recursion.

\documentclass[11pt,a4paper]{article}

%% ---------- packages ----------
\usepackage[utf8]{inputenc}
\usepackage{amsmath,amssymb,amsthm,mathtools}
\usepackage[margin=2.5cm]{geometry}
\usepackage{hyperref}
\usepackage{cleveref}
\usepackage{booktabs}
\usepackage{enumitem}

%% ---------- theorem environments ----------
\newtheorem{theorem}{Theorem}
\newtheorem{lemma}[theorem]{Lemma}
\newtheorem{corollary}[theorem]{Corollary}
\newtheorem{conjecture}[theorem]{Conjecture}
\theoremstyle{definition}
\newtheorem{definition}[theorem]{Definition}
\newtheorem{assumption}{Assumption}
\theoremstyle{remark}
\newtheorem{remark}[theorem]{Remark}

%% ---------- macros ----------
\DeclareMathOperator{\tr}{tr}
\DeclareMathOperator{\spn}{span}
\DeclareMathOperator{\diag}{diag}
\DeclareMathOperator{\KL}{KL}
\newcommand{\R}{\mathbb{R}}
\newcommand{\E}{\mathbb{E}}
\newcommand{\Prob}{\mathbb{P}}
\newcommand{\cM}{\mathcal{M}}
\newcommand{\cS}{\mathcal{S}}
\newcommand{\cX}{\mathcal{X}}
\newcommand{\cD}{\mathcal{D}}
\newcommand{\cC}{\mathcal{C}}
\newcommand{\cF}{\mathcal{F}}
\newcommand{\cG}{\mathcal{G}}
\newcommand{\cH}{\mathcal{H}}
\newcommand{\bJ}{\mathbf{J}}
\newcommand{\norm}[1]{\left\|#1\right\|}

%% ---------- metadata ----------
\title{\Large Discrete Lyapunov Stability of SGD in the Interpolation Regime:\\[4pt]
  \large Per-Sample Analysis and Generalization Bounds}
\author{
  C$\Lambda$ / Lightman\\
  \texttt{Lightman.chang@gmail.com}
}
\date{}

%% =================================================================
\begin{document}
\maketitle

%% =================================================================
%% ABSTRACT
%% =================================================================
\begin{abstract}
We study the implicit regularization of stochastic gradient descent (SGD) in the
interpolation regime, where the training loss reaches zero and the continuous-time
stochastic differential equation (SDE) approximation degenerates because the gradient
noise covariance vanishes at every interpolating solution.  Working directly with the
discrete SGD recursion, we introduce the \emph{per-sample Lyapunov exponent}
$\lambda_k = \frac{1}{n}\ln|1 - \eta h_k|$, where $h_k = \|J_k\|^2$ is the
per-sample curvature and $\eta$ is the learning rate.  Under a Jacobian orthogonality
assumption, we establish that every interpolating solution with
$h_{\max} > 2/\eta$ is almost surely unstable (positive maximal Lyapunov exponent),
while every solution with $h_{\max} < 2/\eta$ is almost surely asymptotically
stable.  We then derive a closed-loop chain:
\emph{SGD stability} $\Rightarrow$ \emph{Jacobian complexity control}
($\cC \leq \sqrt{2/\eta}$) $\Rightarrow$ \emph{Rademacher complexity bound}
$\Rightarrow$ \emph{generalization bound} $O(\rho/\sqrt{\eta n})$.
This mechanism is primarily a discrete phenomenon that does not rely on any
SDE approximation.  We additionally provide a PAC-Bayes bound that eliminates the
dependence on the neighborhood radius~$\rho$.  Our analysis clarifies the distinction
between per-sample stability ($h_{\max} < 2/\eta$) and aggregate spectral conditions
such as $\tr(H) \leq 2/\eta$, and connects stability-based selection to
generalization through an explicit, verifiable complexity measure.
\end{abstract}

%% =================================================================
\section{Introduction}\label{sec:intro}
%% =================================================================

Modern overparameterized neural networks routinely achieve zero training loss---a
regime known as \emph{interpolation}---yet generalize well to unseen data.
Understanding why stochastic gradient descent (SGD) selects, among the continuum of
interpolating solutions, those that generalize is a central open problem in the
theory of deep learning.

\paragraph{Background.}
A fruitful line of work models SGD through its continuous-time SDE approximation
\cite{LiTaiE2017,SmithLe2018}:
\[
  d\theta = -\nabla L(\theta)\,dt + \sqrt{\tfrac{\eta}{B}\,\Sigma(\theta)}\;dW_t,
\]
where $\Sigma(\theta) = \frac{1}{n}\sum_{i=1}^n \nabla\ell_i(\theta)
\nabla\ell_i(\theta)^\top - \nabla L(\theta)\nabla L(\theta)^\top$ is the gradient
noise covariance.  The multiplicative structure of $\Sigma$ biases the
Fokker--Planck stationary distribution toward minima with low noise-to-curvature
ratio $\Sigma/H$ \cite{ChaudhariSoatto2018}.  This provides a satisfying explanation
of implicit regularization---\emph{when} $\Sigma > 0$.

In the interpolation regime, however, every interpolating solution $\theta^*$
satisfies $\nabla\ell_i(\theta^*) = 0$ for all $i$, so $\Sigma(\theta^*) = 0$.
The SDE approximation degenerates and the Fokker--Planck analysis breaks down
entirely.

\paragraph{Known results.}
Wu and Su~\cite{WuSu2023} proved that SGD cannot converge to interpolating solutions
at which $\tr(H) > 2/\eta$, establishing an aggregate spectral condition for
instability.  Wu, Wang, and Su~\cite{WuWangSu2022} studied the sharpness reduction
phenomenon.  Dinh et al.~\cite{Dinh2017} demonstrated that sharpness (Hessian
eigenvalues) alone is not reparameterization-invariant and hence cannot serve as a
sole predictor of generalization.

\paragraph{The gap.}
The condition $\tr(H) > 2/\eta$ is an aggregate condition: it sums over all
eigenvalues and does not distinguish between having one very large per-sample
curvature versus many moderately large ones.  Moreover, the connection from instability
to generalization has remained informal.

\paragraph{Our contribution.}
We provide a complete, self-contained analysis that:
\begin{enumerate}[label=(\roman*)]
  \item introduces the per-sample Lyapunov exponent
    $\lambda_k = \frac{1}{n}\ln|1 - \eta h_k|$ and establishes that stability is
    governed by $h_{\max} < 2/\eta$ rather than the aggregate $\tr(H) < 2/\eta$;
  \item proves that stable solutions have controlled Jacobian complexity
    $\cC(\theta^*) \leq \sqrt{2/\eta}$;
  \item derives a Rademacher-based generalization bound of order
    $O(\rho/\sqrt{\eta n})$ and a PAC-Bayes bound of order
    $O(\sqrt{d\ln(1/\eta)/n})$;
  \item assembles these into a closed-loop theorem: SGD stability $\Rightarrow$
    complexity control $\Rightarrow$ generalization.
\end{enumerate}

The entire analysis works directly on the discrete SGD recursion and does not
invoke any SDE approximation.  The key mechanism---linear stability filtering via
random matrix products---is primarily a discrete phenomenon.

\paragraph{Paper structure.}
\Cref{sec:setup} introduces notation and the interpolation setting.
\Cref{sec:definitions} presents the core definitions (critical learning rate,
good/bad solutions, Jacobian complexity).
\Cref{sec:assumptions} states the assumptions.
\Cref{sec:main} presents the main results: two lemmas, six theorems, and one
corollary.
\Cref{sec:proofs} contains the complete proofs.
\Cref{sec:reparam} discusses reparameterization invariance.
\Cref{sec:discussion} provides connections to the non-interpolation regime,
testable predictions, limitations, and open problems.
\Cref{sec:related} reviews related work.

%% =================================================================
\section{Setup and Notation}\label{sec:setup}
%% =================================================================

Let $S = \{(x_i, y_i)\}_{i=1}^n$ be a training set with
$(x_i, y_i) \stackrel{\text{i.i.d.}}{\sim} \cD$.
Consider a model $f_\theta \colon \cX \to \R$ parameterized by
$\theta \in \R^d$, with $d > n$ (overparameterized regime).
Define the per-sample MSE loss and the empirical risk:
\[
  \ell_i(\theta) = \tfrac{1}{2}\bigl(f_\theta(x_i) - y_i\bigr)^2,
  \qquad
  L(\theta) = \frac{1}{n}\sum_{i=1}^n \ell_i(\theta).
\]

\paragraph{Interpolation manifold.}
The set of global minimizers of $L$ with zero training loss is
\[
  \cM = \bigl\{\theta \in \R^d : f_\theta(x_i) = y_i,\;\forall\, i \in [n]\bigr\}.
\]
Since $d > n$, generically $\cM$ is a smooth submanifold of dimension at least $d - n$.

\paragraph{Per-sample quantities at $\theta^* \in \cM$.}
\begin{itemize}[leftmargin=2em]
  \item \textbf{Jacobian vectors:} $J_i = \nabla_\theta f_{\theta^*}(x_i) \in \R^d$.
  \item \textbf{Jacobian matrix:} $\bJ \in \R^{n \times d}$, with the $i$-th row
    equal to $J_i^\top$.
  \item \textbf{Per-sample Hessian:}
    $H_i = J_i J_i^\top$.
    At $\theta^* \in \cM$ the residual $r_i(\theta^*) = f_{\theta^*}(x_i) - y_i = 0$,
    so $\nabla^2 \ell_i(\theta^*) = H_i$.
  \item \textbf{Per-sample curvature:} $h_i = \|J_i\|^2$.
  \item \textbf{Maximum curvature:} $h_{\max} = \max_{i \in [n]} h_i$.
  \item \textbf{Active subspace:} $\cS = \spn(J_1, \ldots, J_n)$.
\end{itemize}

\paragraph{SGD recursion.}
\[
  \theta_{t+1} = \theta_t - \eta\,\nabla\ell_{z_t}(\theta_t),
  \qquad z_t \sim \mathrm{Uniform}([n]) \;\text{i.i.d.}
\]

\paragraph{Population risk.}
$L_{\mathrm{pop}}(\theta) = \E_{(x,y)\sim\cD}[\ell(f_\theta(x), y)]$.

%% =================================================================
\section{Definitions}\label{sec:definitions}
%% =================================================================

\begin{definition}[Critical learning rate]\label{def:eta-c}
  For an interpolating solution $\theta^* \in \cM$, the \emph{critical learning rate}
  is
  \[
    \eta_c(\theta^*) = \frac{2}{h_{\max}(\theta^*)}.
  \]
\end{definition}

\begin{definition}[Good and bad solutions]\label{def:good-bad}
  Given a learning rate $\eta > 0$:
  \begin{itemize}
    \item $\theta^*$ is \textbf{$\eta$-stable} (good, flat) if
      $\eta < \eta_c(\theta^*)$, equivalently $\eta\, h_{\max} < 2$.
    \item $\theta^*$ is \textbf{$\eta$-unstable} (bad, sharp) if
      $\eta > \eta_c(\theta^*)$, equivalently $\eta\, h_{\max} > 2$.
  \end{itemize}
  The set of $\eta$-stable interpolating solutions is
  $\cM_\eta = \{\theta^* \in \cM : h_{\max}(\theta^*) < 2/\eta\}$.
\end{definition}

\begin{definition}[Jacobian complexity]\label{def:jac-complex}
  \[
    \cC(\theta^*) = \sqrt{\frac{1}{n}\sum_{i=1}^n h_i(\theta^*)}
    = \frac{\|\bJ(\theta^*)\|_F}{\sqrt{n}}.
  \]
\end{definition}

\begin{remark}[Physical intuition]\label{rem:intuition}
  The quantity $h_i = \|\nabla_\theta f_{\theta^*}(x_i)\|^2$ measures the sensitivity
  of the model output at $x_i$ with respect to the parameters.  A good interpolating
  solution is one where no training sample requires the model to be in a
  high-sensitivity configuration.  A bad interpolating solution has at least one
  sample $x_i$ for which the model must precisely ``align'' its parameters to
  interpolate; such alignment cannot be maintained under the stochastic perturbations
  of SGD.
\end{remark}

%% =================================================================
\section{Assumptions}\label{sec:assumptions}
%% =================================================================

\begin{assumption}[Local smoothness]\label{A:smooth}
  $f_\theta(x)$ is twice continuously differentiable in $\theta$.  In a neighborhood
  $B(\theta^*, r_0)$ of $\theta^*$,
  \[
    \|\nabla^2_\theta f_\theta(x_i)\|_{\mathrm{op}} \leq B_2
    \quad \text{for all } i \in [n].
  \]
\end{assumption}

\begin{assumption}[Jacobian independence]\label{A:indep}
  $\{J_1, \ldots, J_n\}$ are linearly independent, so $\dim \cS = n$.
\end{assumption}

\begin{assumption}[Jacobian orthogonality]\label{A:orth}
  $J_i^\top J_j = 0$ for all $i \neq j$.
\end{assumption}

\begin{remark}[On the orthogonality assumption]\label{rem:orth}
  \Cref{A:orth} is a technical assumption that enables complete decoupling of the
  stability analysis into $n$ independent scalar problems.  We emphasize that this
  is a genuine restriction: whitening the kernel matrix $K = \bJ\bJ^\top$
  diagonalizes $K$ but does \emph{not} make the Jacobian vectors $J_i$ orthogonal
  in parameter space (it only orthogonalizes them in the $n$-dimensional output
  space).  Thus the NTK regime does not automatically satisfy \Cref{A:orth}.
  The qualitative conclusion---$\eta h_{\max} > 2$ implies instability---is
  expected to hold in the general (non-orthogonal) case by Furstenberg's theory
  of random matrix products \cite{furstenberg1960products}, but rigorous
  quantitative bounds in the non-commutative setting remain open
  (see \Cref{sec:discussion}).  We present fully rigorous proofs under
  \Cref{A:orth}, which has the virtue of yielding exact, closed-form Lyapunov
  exponents.
\end{remark}

%% =================================================================
\section{Main Results}\label{sec:main}
%% =================================================================

We state the results in logical order.  All proofs are deferred to \Cref{sec:proofs}.

\begin{lemma}[Linearization of SGD near an interpolating solution]\label{lem:lin}
  Let $\theta^* \in \cM$ and suppose \Cref{A:smooth} holds.
  Define $\delta_t = \theta_t - \theta^*$.  When $\|\delta_t\| < r_0$,
  \[
    \delta_{t+1} = A_{z_t}\,\delta_t + R_{z_t}(\delta_t),
  \]
  where $A_i = I - \eta\,J_i J_i^\top$ and
  $\|R_i(\delta)\| \leq C\eta\|\delta\|^2$ for a constant $C > 0$.
\end{lemma}

\begin{lemma}[Spectrum of the rank-one update]\label{lem:rank1}
  The matrix $A_i = I - \eta J_i J_i^\top$ has eigenvalues:
  \begin{itemize}
    \item $1 - \eta h_i$, with eigenvector $J_i$;
    \item $1$ with multiplicity $d-1$, corresponding to the eigenspace $J_i^\perp$.
  \end{itemize}
  The operator norm satisfies
  $\|A_i\|_{\mathrm{op}} = \max(1, |1 - \eta h_i|)$.
  In particular, $|1 - \eta h_i| \leq 1$ if and only if $0 \leq \eta h_i \leq 2$.
\end{lemma}

\begin{theorem}[Instability of bad interpolating solutions]\label{thm:instability}
  Let $\theta^* \in \cM$ and suppose \Cref{A:smooth}--\Cref{A:orth} hold, with
  $\eta\, h_{\max} > 2$.  Then the maximal Lyapunov exponent of the linearized SGD
  at $\theta^*$ is
  \[
    \lambda(\eta, \theta^*) = \max_{k \in [n]} \frac{1}{n}\ln|1 - \eta h_k| > 0.
  \]
  Consequently, $\theta^*$ is \textbf{almost surely unstable} under SGD: for almost
  every initial perturbation $\delta_0 \neq 0$ with nonzero component in $\cS$, the
  linearized system satisfies $\|\delta_t\| \to \infty$.
\end{theorem}

\begin{theorem}[Stability of good interpolating solutions]\label{thm:stability}
  Let $\theta^* \in \cM$ and suppose \Cref{A:smooth}--\Cref{A:orth} hold, with
  $\eta\, h_{\max} < 2$.  Then all Lyapunov exponents are negative:
  \[
    \lambda_k = \frac{1}{n}\ln|1 - \eta h_k| < 0,
    \quad \forall\, k \in [n].
  \]
  The solution $\theta^*$ is \textbf{almost surely asymptotically stable} under
  linearized SGD: the active-subspace component decays exponentially,
  \[
    \|\delta_t|_{\cS}\| \leq \|\delta_0|_{\cS}\| \cdot
    \exp\!\Bigl(\max_k \lambda_k \cdot t + o(t)\Bigr),
  \]
  while $P_{\cS^\perp}\delta_t = P_{\cS^\perp}\delta_0$ remains constant.
\end{theorem}

\begin{theorem}[Jacobian complexity controls Rademacher complexity]\label{thm:rad}
  Let $\theta^* \in \cM$ and suppose \Cref{A:smooth} holds.  Define
  \[
    \cG_\rho = \bigl\{x \mapsto f_\theta(x) - f_{\theta^*}(x)
    : \theta \in B(\theta^*, \rho)\bigr\}.
  \]
  Then the empirical Rademacher complexity satisfies
  \begin{equation}\label{eq:rad}
    \hat{\mathcal{R}}_S(\cG_\rho)
    \leq \frac{\rho\,\cC(\theta^*)}{\sqrt{n}} + \frac{B_2\,\rho^2}{2}.
  \end{equation}
\end{theorem}

\begin{theorem}[Generalization bound]\label{thm:gen}
  Let $\theta^* \in \cM$ and suppose \Cref{A:smooth} holds.
  Suppose the loss $\ell(\hat{y}, y)$ is $L_\ell$-Lipschitz in its first argument
  and takes values in $[0, \bar{\ell}]$.
  Then with probability at least $1 - \delta$ over the draw of $S$:
  \begin{equation}\label{eq:gen}
    L_{\mathrm{pop}}(\theta^*)
    \leq \underbrace{L(\theta^*)}_{=\,0}
    + \frac{2 L_\ell\,\rho\,\cC(\theta^*)}{\sqrt{n}}
    + L_\ell\, B_2\,\rho^2
    + \bar{\ell}\sqrt{\frac{\ln(2/\delta)}{2n}}.
  \end{equation}
\end{theorem}

\begin{theorem}[Closed-loop implicit regularization---Main Theorem]\label{thm:main}
  Let \Cref{A:smooth}--\Cref{A:orth} hold and $\cM \neq \varnothing$.
  Define $\cM_\eta = \{\theta^* \in \cM : h_{\max}(\theta^*) < 2/\eta\}$.
  Consider SGD with learning rate $\eta$.
  \begin{enumerate}[label=\textbf{(\Roman*)}]
    \item \textbf{Selection.}\;
      For every interpolating solution in $\cM \setminus \cM_\eta$
      (i.e., with $\eta h_{\max} > 2$),
      the maximal Lyapunov exponent of the \emph{linearized} SGD satisfies
      $\lambda = \frac{1}{n}\ln(\eta h_{\max} - 1) > 0$,
      so these solutions are almost surely unstable under linearized SGD.
      Conversely, for solutions in $\cM_\eta$ (with $\eta h_{\max} < 2$),
      all Lyapunov exponents in $\cS$ are negative and the $\cS$-direction
      perturbation decays exponentially under linearized SGD.
      The boundary case $\eta h_{\max} = 2$ is degenerate
      ($\lambda = 0$) and is excluded from both sets.

    \item \textbf{Complexity control.}\;
      Every solution in $\cM_\eta$ satisfies $h_i \leq 2/\eta$ for all $i$, hence
      \[
        \cC(\theta^*)^2 = \frac{1}{n}\sum_{i=1}^n h_i
        \leq h_{\max} < \frac{2}{\eta}.
      \]

    \item \textbf{Generalization.}\;
      With probability at least $1 - \delta$, any SGD-selected solution
      $\theta^* \in \cM_\eta$ satisfies
      \[
        L_{\mathrm{pop}}(\theta^*)
        \leq \frac{2\sqrt{2}\, L_\ell\,\rho}{\sqrt{\eta\, n}}
        + L_\ell\, B_2\,\rho^2
        + \bar{\ell}\sqrt{\frac{\ln(2/\delta)}{2n}}.
      \]

    \item \textbf{Monotonicity.}\;
      $\eta_1 < \eta_2 \;\Rightarrow\; \cM_{\eta_2} \subseteq \cM_{\eta_1}
      \;\Rightarrow\; \sup_{\cM_{\eta_2}} \cC \leq \sup_{\cM_{\eta_1}} \cC$.
      Increasing $\eta$ tightens the generalization bound.
  \end{enumerate}
\end{theorem}

\begin{theorem}[PAC-Bayes bound]\label{thm:pb}
  Let $\theta^* \in \cM$ and suppose \Cref{A:smooth} holds.
  Let the prior be $P = \mathcal{N}(0, \sigma_0^2 I_d)$ and the posterior be
  $Q = \mathcal{N}(\theta^*, \sigma^2 I_d)$.
  Suppose $\ell(\hat{y}, y) = \frac{1}{2}(\hat{y} - y)^2 \in [0, \bar{\ell}]$.
  Then with probability at least $1 - \delta$:
  \begin{equation}\label{eq:pb}
    \E_{\theta \sim Q}\bigl[L_{\mathrm{pop}}(\theta)\bigr]
    \leq \frac{\sigma^2\,\cC(\theta^*)^2}{2} + O(\sigma^4)
    + \sqrt{
      \frac{
        \frac{d}{2}\!\left(\frac{\sigma^2}{\sigma_0^2} - 1
        - \ln\frac{\sigma^2}{\sigma_0^2}\right)
        + \frac{\|\theta^*\|^2}{2\sigma_0^2}
        + \ln\frac{2\sqrt{n}}{\delta}
      }{2n}
    }.
  \end{equation}
\end{theorem}

\begin{corollary}[PAC-Bayes bound for SGD-selected solutions]\label{cor:pb}
  Suppose SGD selects $\theta^* \in \cM_\eta$, so that
  $\cC(\theta^*)^2 \leq 2/\eta$.  Set $\sigma^2 = \alpha\eta$ for a constant
  $\alpha > 0$.  Then:
  \[
    \E_{\theta \sim Q}[L_{\mathrm{pop}}(\theta)]
    \leq \alpha + O(\eta^2)
    + \sqrt{
      \frac{
        \frac{d}{2}\!\left(\frac{\alpha\eta}{\sigma_0^2} - 1
        - \ln\frac{\alpha\eta}{\sigma_0^2}\right)
        + \frac{\|\theta^*\|^2}{2\sigma_0^2}
        + \ln\frac{2\sqrt{n}}{\delta}
      }{2n}
    }.
  \]
  In particular, when $\alpha\eta \ll \sigma_0^2$ the KL term is dominated by
  $\frac{d}{2}\ln\frac{\sigma_0^2}{\alpha\eta}$, and the bound scales as
  $O\bigl(\sqrt{d\,\ln(1/\eta)\,/\,n}\bigr)$.
\end{corollary}

%% =================================================================
\section{Proofs}\label{sec:proofs}
%% =================================================================

\subsection{Proof of Lemma~\ref{lem:lin} (Linearization)}

\begin{proof}
\textbf{Step~1 (Residual expansion).}\;
At $\theta^*$, the residual is $r_i(\theta^*) = f_{\theta^*}(x_i) - y_i = 0$
by the interpolation condition.  By Taylor expansion around $\theta^*$:
\[
  r_i(\theta^* + \delta)
  = \underbrace{r_i(\theta^*)}_{=\,0} + J_i^\top \delta + O(\|\delta\|^2)
  = J_i^\top \delta + O(\|\delta\|^2).
\]

\textbf{Step~2 (Jacobian expansion).}\;
Similarly,
\[
  \nabla_\theta f_{\theta^* + \delta}(x_i)
  = J_i + \nabla^2_\theta f_{\theta^*}(x_i)\,\delta + O(\|\delta\|^2)
  = J_i + O(\|\delta\|).
\]

\textbf{Step~3 (Loss gradient expansion).}\;
Since $\nabla\ell_i(\theta) = r_i(\theta) \cdot \nabla_\theta f_\theta(x_i)$,
we have
\begin{align*}
  \nabla\ell_i(\theta^* + \delta)
  &= \bigl(J_i^\top \delta + O(\|\delta\|^2)\bigr)
     \bigl(J_i + O(\|\delta\|)\bigr) \\
  &= J_i(J_i^\top \delta) + O(\|\delta\|^2) \\
  &= J_i J_i^\top \delta + O(\|\delta\|^2) \\
  &= H_i\,\delta + O(\|\delta\|^2).
\end{align*}
For the $O(\|\delta\|^2)$ remainder: the cross terms are
$(J_i^\top \delta) \cdot \nabla^2_\theta f_{\theta^*}(x_i)\,\delta$
(of magnitude $\sqrt{h_i}\,\|\delta\| \cdot B_2\,\|\delta\|$)
and
$\frac{1}{2}(\delta^\top \nabla^2_\theta f_{\theta^*}(x_i)\,\delta) \cdot J_i$
(of magnitude $\frac{B_2}{2}\|\delta\|^2\sqrt{h_i}$),
plus higher-order terms.  Therefore
$\|\nabla\ell_i(\theta^* + \delta) - H_i\,\delta\| \leq C_1\|\delta\|^2$
where $C_1 = B_2\bigl(\sqrt{h_{\max}} + \frac{B_2}{2}\sqrt{h_{\max}}\bigr)$.

\textbf{Step~4 (SGD update).}\;
\begin{align*}
  \delta_{t+1}
  &= \delta_t - \eta\,\nabla\ell_{z_t}(\theta^* + \delta_t) \\
  &= \delta_t - \eta\,H_{z_t}\,\delta_t - \eta \cdot O(\|\delta_t\|^2) \\
  &= (I - \eta\,H_{z_t})\,\delta_t + R_{z_t}(\delta_t),
\end{align*}
where $\|R_i(\delta)\| \leq \eta\, C_1\,\|\delta\|^2$.  Setting $C = C_1$
completes the proof.
\end{proof}

\subsection{Proof of Lemma~\ref{lem:rank1} (Rank-one spectrum)}

\begin{proof}
\textbf{Action on $J_i$:}\;
$A_i\, J_i = (I - \eta\, J_i J_i^\top)\, J_i
= J_i - \eta\,\|J_i\|^2\, J_i = (1 - \eta h_i)\, J_i$.

\textbf{Action on $u \perp J_i$:}\;
$A_i\, u = (I - \eta\, J_i J_i^\top)\, u = u - \eta\,(J_i^\top u)\, J_i = u$.

Thus the eigenvalues are $1 - \eta h_i$ (eigenvector $J_i$) and $1$ with
multiplicity $d - 1$ (eigenspace $J_i^\perp$).

The operator norm is
$\|A_i\|_{\mathrm{op}} = \max(|1 - \eta h_i|, 1)$.

Finally,
$|1 - \eta h_i| \leq 1
\;\Leftrightarrow\; -1 \leq 1 - \eta h_i \leq 1
\;\Leftrightarrow\; 0 \leq \eta h_i \leq 2$.
\end{proof}

\subsection{Proof of Theorem~\ref{thm:instability} (Instability)}

\begin{proof}
\textbf{Step~1 (Coordinate decoupling under orthogonality).}\;
By \Cref{A:orth}, $J_i^\top J_j = 0$ for $i \neq j$.
Construct an orthonormal basis of $\cS$:
\[
  e_k = \frac{J_k}{\sqrt{h_k}}, \quad k = 1, \ldots, n.
\]
Then $e_k^\top e_l = \frac{J_k^\top J_l}{\sqrt{h_k\, h_l}} = \delta_{kl}$.

\textbf{Step~2 (Matrix representation of $A_i$ on $\cS$).}\;
The $(k, l)$-entry of $A_i$ restricted to $\cS$ in the basis $\{e_k\}$ is
\begin{align*}
  (\hat{A}_i)_{kl}
  &= e_k^\top A_i\, e_l
  = e_k^\top(I - \eta\, J_i J_i^\top)\, e_l
  = \delta_{kl} - \eta\,(e_k^\top J_i)(J_i^\top e_l).
\end{align*}
By orthogonality, $e_k^\top J_i = \frac{J_k^\top J_i}{\sqrt{h_k}}
= \sqrt{h_i}\,\delta_{ki}$.  Therefore
\[
  (\hat{A}_i)_{kl} = \delta_{kl} - \eta\, h_i\,\delta_{ki}\,\delta_{li}.
\]
This means $\hat{A}_i$ is a \textbf{diagonal matrix}:
\[
  (\hat{A}_i)_{kk} =
  \begin{cases}
    1 - \eta h_i, & k = i, \\
    1, & k \neq i.
  \end{cases}
\]

\textbf{Step~3 (Coordinate evolution).}\;
Let $\xi_t \in \R^n$ denote the coordinates of $\delta_t|_{\cS}$ in the basis
$\{e_k\}$: $(\xi_t)_k = e_k^\top \delta_t$.  By the diagonal structure of
$\hat{A}_i$:
\[
  (\xi_{t+1})_k = (\hat{A}_{z_t})_{kk} \cdot (\xi_t)_k
  = \begin{cases}
    (1 - \eta h_k)(\xi_t)_k, & z_t = k, \\
    (\xi_t)_k, & z_t \neq k.
  \end{cases}
\]
Each coordinate $k$ evolves \textbf{independently}: there is no coupling between
different coordinates.

\textbf{Step~4 (Product representation).}\;
\[
  (\xi_T)_k = (\xi_0)_k \cdot \prod_{t=1}^{T} (\hat{A}_{z_t})_{kk}.
\]
Taking logarithms:
\[
  \ln|(\xi_T)_k| = \ln|(\xi_0)_k| + \sum_{t=1}^{T} X_t^{(k)},
\]
where $X_t^{(k)} = \ln|(\hat{A}_{z_t})_{kk}|$.

\textbf{Step~5 (Distribution of $X_t^{(k)}$).}\;
Since $z_t \sim \mathrm{Uniform}([n])$:
\[
  X_t^{(k)} =
  \begin{cases}
    \ln|1 - \eta h_k|, & \text{with probability } 1/n, \\
    0, & \text{with probability } (n-1)/n.
  \end{cases}
\]
The mean is
$\E[X_t^{(k)}] = \frac{1}{n}\ln|1 - \eta h_k| \triangleq \lambda_k$.

The variance is
$\mathrm{Var}(X_t^{(k)})
= \frac{1}{n}(\ln|1 - \eta h_k|)^2 - \frac{1}{n^2}(\ln|1 - \eta h_k|)^2
= \frac{n-1}{n^2}(\ln|1 - \eta h_k|)^2 < \infty$.

\textbf{Step~6 (Strong law of large numbers).}\;
The sequence $\{X_t^{(k)}\}_{t \geq 1}$ is i.i.d.\ (since $z_1, z_2, \ldots$
are i.i.d.) with finite mean and variance.
By the Kolmogorov strong law of large numbers:
\[
  \frac{1}{T}\sum_{t=1}^{T} X_t^{(k)}
  \xrightarrow{T \to \infty} \lambda_k
  = \frac{1}{n}\ln|1 - \eta h_k|
  \quad \text{almost surely.}
\]
Hence
\[
  \frac{1}{T}\ln|(\xi_T)_k| \to \lambda_k
  \quad \text{almost surely.}
\]

\textbf{Step~7 (Sign of the Lyapunov exponent).}\;

\emph{Case 1:} $\eta h_k < 2$.
Then $|1 - \eta h_k| < 1$, so $\ln|1 - \eta h_k| < 0$ and $\lambda_k < 0$.

\emph{Case 2:} $\eta h_k > 2$.
Then $|1 - \eta h_k| = \eta h_k - 1 > 1$, so $\ln|1 - \eta h_k| > 0$ and
$\lambda_k > 0$.

\textbf{Step~8 (Instability conclusion).}\;
Let $k^* = \arg\max_k h_k$.  By hypothesis,
$\eta h_{k^*} = \eta h_{\max} > 2$, so $\lambda_{k^*} > 0$.
Provided $(\xi_0)_{k^*} \neq 0$ (i.e., the initial perturbation has a nonzero
component in the $J_{k^*}$ direction),
\[
  |(\xi_T)_{k^*}|
  = |(\xi_0)_{k^*}| \cdot \exp\!\bigl(\lambda_{k^*} \cdot T + o(T)\bigr)
  \to \infty
  \quad \text{almost surely.}
\]
Therefore $\|\delta_T\| \geq |(\xi_T)_{k^*}| \to \infty$, and $\theta^*$ is
almost surely unstable under the linearized SGD.

\textbf{Step~9 (Maximal Lyapunov exponent).}\;
The maximal Lyapunov exponent of the system on $\cS$ is
\[
  \lambda = \max_{k \in [n]} \lambda_k
  = \max_k \frac{1}{n}\ln|1 - \eta h_k|
  \geq \lambda_{k^*}
  = \frac{1}{n}\ln(\eta h_{\max} - 1) > 0. \qedhere
\]
\end{proof}

\subsection{Proof of Theorem~\ref{thm:stability} (Stability)}

\begin{proof}
All steps follow the same structure as the proof of \Cref{thm:instability}, with the
sign reversed.

\textbf{Steps 1--6:}\;
Identical to Steps~1--6 in the proof of \Cref{thm:instability}.  We obtain the
per-coordinate Lyapunov exponent
$\lambda_k = \frac{1}{n}\ln|1 - \eta h_k|$ for each $k \in [n]$.

\textbf{Step~7 (All exponents are negative).}\;
By hypothesis, $\eta h_k \leq \eta h_{\max} < 2$ for all $k$.  By \Cref{lem:rank1},
$|1 - \eta h_k| < 1$ for every $k$, so $\lambda_k < 0$ for all $k \in [n]$.

\textbf{Step~8 (Asymptotic stability).}\;
Since $\frac{1}{T}\ln|(\xi_T)_k| \to \lambda_k < 0$ almost surely (by Step~6),
we have $|(\xi_T)_k| \to 0$ for every $k$.  Therefore
\[
  \|\delta_T|_{\cS}\| = \sqrt{\sum_{k=1}^n (\xi_T)_k^2} \to 0
  \quad \text{almost surely.}
\]
The decay rate is governed by the slowest coordinate:
$\lambda_{\max} = \max_k \lambda_k < 0$, so
\[
  \|\delta_T|_{\cS}\| \leq \|\delta_0|_{\cS}\| \cdot
  \exp(\lambda_{\max}\, T + o(T)).
\]

\textbf{Behavior in $\cS^\perp$.}\;
For any $v \in \cS^\perp$, $A_i\, v = v$ by \Cref{lem:rank1} (since
$v \perp J_i$ for all $i$).  Therefore
$P_{\cS^\perp}\delta_t = P_{\cS^\perp}\delta_0$, which remains constant.
\end{proof}

\subsection{Proof of Theorem~\ref{thm:rad} (Rademacher bound)}

\begin{proof}
\textbf{Step~1 (Taylor expansion).}\;
For $\theta = \theta^* + \delta$ with $\|\delta\| \leq \rho$:
\[
  f_\theta(x_i) - f_{\theta^*}(x_i) = J_i^\top \delta + r_i(\delta),
\]
where $r_i(\delta) = \frac{1}{2}\delta^\top \nabla^2_\theta f_{\tilde{\theta}}(x_i)
\,\delta$ for some $\tilde{\theta}$ on the segment $[\theta^*, \theta]$, by Taylor's
theorem with Lagrange remainder.  By \Cref{A:smooth},
$|r_i(\delta)| \leq \frac{B_2}{2}\|\delta\|^2 \leq \frac{B_2\,\rho^2}{2}$.

\textbf{Step~2 (Rademacher decomposition).}\;
Let $\sigma_1, \ldots, \sigma_n$ be i.i.d.\ Rademacher random variables.
\begin{align*}
  \hat{\mathcal{R}}_S(\cG_\rho)
  &= \E_\sigma\!\left[
    \sup_{\|\delta\| \leq \rho} \frac{1}{n}\sum_{i=1}^n
    \sigma_i\bigl(J_i^\top \delta + r_i(\delta)\bigr)
  \right] \\
  &\leq \E_\sigma\!\left[
    \sup_{\|\delta\| \leq \rho} \frac{1}{n}\sum_{i=1}^n \sigma_i\, J_i^\top \delta
  \right]
  + \frac{B_2\,\rho^2}{2}.
\end{align*}
The second inequality holds because
$\sup_{\|\delta\| \leq \rho} \frac{1}{n}\sum_i \sigma_i\bigl(J_i^\top \delta
+ r_i(\delta)\bigr)
\leq \sup_{\|\delta\| \leq \rho} \frac{1}{n}\sum_i \sigma_i\, J_i^\top \delta
+ \sup_{\|\delta\| \leq \rho} \frac{1}{n}\sum_i |\sigma_i\, r_i(\delta)|
\leq \sup_{\|\delta\| \leq \rho} \frac{1}{n}\sum_i \sigma_i\, J_i^\top \delta
+ \frac{B_2\,\rho^2}{2}$,
where the last term is deterministic since $|r_i(\delta)| \leq \frac{B_2\rho^2}{2}$
regardless of $\sigma$.

\textbf{Step~3 (Linear part computation).}\;
\begin{align*}
  \sup_{\|\delta\| \leq \rho}
  \frac{1}{n}\sum_{i=1}^n \sigma_i\, J_i^\top \delta
  &= \sup_{\|\delta\| \leq \rho}
     \frac{1}{n}\!\left(\sum_{i=1}^n \sigma_i\, J_i\right)^{\!\top} \delta
  = \frac{\rho}{n}\left\|\sum_{i=1}^n \sigma_i\, J_i\right\|.
\end{align*}
The last equality follows from $\sup_{\|\delta\| \leq \rho} v^\top \delta
= \rho\,\|v\|$, attained at $\delta^* = \rho\, v/\|v\|$.

\textbf{Step~4 (Expectation estimate via Jensen's inequality).}\;
Since $\sqrt{\cdot}$ is concave, by Jensen's inequality:
\[
  \E_\sigma\!\left[\left\|\sum_{i=1}^n \sigma_i\, J_i\right\|\right]
  \leq \sqrt{\E_\sigma\!\left[\left\|\sum_{i=1}^n \sigma_i\, J_i\right\|^2\right]}.
\]
Expanding the squared norm:
\begin{align*}
  \E_\sigma\!\left[\left\|\sum_{i=1}^n \sigma_i\, J_i\right\|^2\right]
  &= \E_\sigma\!\left[\sum_{i=1}^n \sum_{j=1}^n
     \sigma_i\, \sigma_j\, J_i^\top J_j\right]
  = \sum_{i=1}^n \sum_{j=1}^n \E[\sigma_i\, \sigma_j] \cdot J_i^\top J_j.
\end{align*}
Since $\sigma_i$ are i.i.d.\ Rademacher:
$\E[\sigma_i\,\sigma_j] = \delta_{ij}$
(because $\E[\sigma_i^2] = 1$ and
$\E[\sigma_i\,\sigma_j] = \E[\sigma_i]\,\E[\sigma_j] = 0$ for $i \neq j$).
Therefore
\[
  \E_\sigma\!\left[\left\|\sum_{i=1}^n \sigma_i\, J_i\right\|^2\right]
  = \sum_{i=1}^n \|J_i\|^2 = \sum_{i=1}^n h_i = n\,\cC(\theta^*)^2.
\]
Hence $\E_\sigma\!\bigl[\|\sum_i \sigma_i\, J_i\|\bigr]
\leq \sqrt{n}\,\cC(\theta^*)$.

\textbf{Step~5 (Combining).}\;
\[
  \hat{\mathcal{R}}_S(\cG_\rho)
  \leq \frac{\rho}{n} \cdot \sqrt{n}\,\cC(\theta^*)
  + \frac{B_2\,\rho^2}{2}
  = \frac{\rho\,\cC(\theta^*)}{\sqrt{n}} + \frac{B_2\,\rho^2}{2}.
  \qedhere
\]
\end{proof}

\subsection{Proof of Theorem~\ref{thm:gen} (Generalization bound)}

\begin{proof}
\textbf{Step~1 (Zero training loss).}\;
Since $\theta^* \in \cM$, we have $f_{\theta^*}(x_i) = y_i$ for all $i$,
so $L(\theta^*) = 0$.

\textbf{Step~2 (Standard Rademacher generalization bound).}\;
Define the function class
$\cH = \{(x,y) \mapsto \ell(f_\theta(x), y) : \theta \in B(\theta^*, \rho)\}$.
By the standard uniform convergence result (Mohri, Rostamizadeh, and Talwalkar,
Theorem~3.1), with probability at least $1 - \delta$ over the draw of $S$:
\[
  \sup_{\theta \in B(\theta^*, \rho)}
  \bigl|L_{\mathrm{pop}}(\theta) - L(\theta)\bigr|
  \leq 2\,\hat{\mathcal{R}}_S(\cH)
  + \bar{\ell}\sqrt{\frac{\ln(2/\delta)}{2n}}.
\]

\textbf{Step~3 (Lipschitz contraction).}\;
By the Ledoux--Talagrand contraction lemma, the $L_\ell$-Lipschitz property of
$\ell$ in its first argument implies
\[
  \hat{\mathcal{R}}_S(\cH) \leq L_\ell\,\hat{\mathcal{R}}_S(\cF_\rho),
\]
where $\cF_\rho = \{x \mapsto f_\theta(x) : \theta \in B(\theta^*, \rho)\}$.
Adding and subtracting the constant $f_{\theta^*}$ does not affect the Rademacher
complexity (since $\E_\sigma[\sum_i \sigma_i\, c] = 0$ for any constant $c$),
so $\hat{\mathcal{R}}_S(\cF_\rho) = \hat{\mathcal{R}}_S(\cG_\rho)$.

\textbf{Step~4 (Substituting \Cref{thm:rad}).}\;
\begin{align*}
  L_{\mathrm{pop}}(\theta^*)
  &\leq L(\theta^*)
  + 2\,L_\ell\,\hat{\mathcal{R}}_S(\cG_\rho)
  + \bar{\ell}\sqrt{\frac{\ln(2/\delta)}{2n}} \\
  &\leq 0
  + 2\,L_\ell\!\left(
    \frac{\rho\,\cC(\theta^*)}{\sqrt{n}} + \frac{B_2\,\rho^2}{2}
  \right)
  + \bar{\ell}\sqrt{\frac{\ln(2/\delta)}{2n}} \\
  &= \frac{2\,L_\ell\,\rho\,\cC(\theta^*)}{\sqrt{n}}
  + L_\ell\, B_2\,\rho^2
  + \bar{\ell}\sqrt{\frac{\ln(2/\delta)}{2n}}.
  \qedhere
\end{align*}
\end{proof}

\subsection{Proof of Theorem~\ref{thm:main} (Main Theorem---Closed-loop implicit regularization)}

\begin{proof}
\textbf{Part (I): Selection.}\;
This follows directly from \Cref{thm:instability} and \Cref{thm:stability}.
For $\theta^* \in \cM \setminus \cM_\eta$, we have $h_{\max}(\theta^*) \geq 2/\eta$,
so $\eta h_{\max} \geq 2$.  When $\eta h_{\max} > 2$, by \Cref{thm:instability}
the maximal Lyapunov exponent is
$\lambda \geq \frac{1}{n}\ln(\eta h_{\max} - 1) > 0$,
and SGD escapes almost surely.
For $\theta^* \in \cM_\eta$, we have $\eta h_{\max} < 2$, and by
\Cref{thm:stability} all Lyapunov exponents are negative and the $\cS$-direction
perturbation decays exponentially.

\textbf{Part (II): Complexity control.}\;
For $\theta^* \in \cM_\eta$, by definition $h_i \leq h_{\max} < 2/\eta$ for
every $i \in [n]$.  Therefore
\[
  \cC(\theta^*)^2 = \frac{1}{n}\sum_{i=1}^n h_i
  \leq h_{\max} < \frac{2}{\eta}.
\]

\textbf{Part (III): Generalization.}\;
Substituting $\cC(\theta^*) \leq \sqrt{2/\eta}$ into \eqref{eq:gen}:
\[
  L_{\mathrm{pop}}(\theta^*)
  \leq \frac{2\,L_\ell\,\rho\,\sqrt{2/\eta}}{\sqrt{n}}
  + L_\ell\, B_2\,\rho^2
  + \bar{\ell}\sqrt{\frac{\ln(2/\delta)}{2n}}
  = \frac{2\sqrt{2}\,L_\ell\,\rho}{\sqrt{\eta\, n}}
  + L_\ell\, B_2\,\rho^2
  + \bar{\ell}\sqrt{\frac{\ln(2/\delta)}{2n}}.
\]

\textbf{Part (IV): Monotonicity.}\;
By definition, $\cM_\eta = \{\theta^* \in \cM : h_{\max} < 2/\eta\}$.
When $\eta$ increases, $2/\eta$ decreases, making the threshold stricter.
Hence $\eta_1 < \eta_2$ implies $\cM_{\eta_2} \subseteq \cM_{\eta_1}$.
Consequently,
$\sup_{\theta^* \in \cM_{\eta_2}} \cC(\theta^*)
\leq \sup_{\theta^* \in \cM_{\eta_1}} \cC(\theta^*)$.
Moreover, $\sup_{\cM_\eta} \cC \leq \sqrt{2/\eta}$ is monotonically decreasing
in $\eta$.
\end{proof}

\subsection{Proof of Theorem~\ref{thm:pb} (PAC-Bayes bound)}

\begin{proof}
\textbf{Step~1 (PAC-Bayes inequality).}\;
By the standard PAC-Bayes theorem (McAllester 1999; Catoni 2007), with
probability at least $1 - \delta$ over the draw of $S$:
\[
  \E_{\theta \sim Q}[L_{\mathrm{pop}}(\theta)]
  \leq \E_{\theta \sim Q}[L(\theta)]
  + \sqrt{\frac{\KL(Q \| P) + \ln(2\sqrt{n}/\delta)}{2n}}.
\]

\textbf{Step~2 (Computing $\E_Q[L]$).}\;
Write $\theta = \theta^* + \epsilon$ where $\epsilon \sim \mathcal{N}(0, \sigma^2 I)$.
By Taylor expansion at $\theta^*$:
\[
  f_\theta(x_i) - y_i
  = J_i^\top \epsilon
  + \tfrac{1}{2}\epsilon^\top \nabla^2_\theta f_{\theta^*}(x_i)\,\epsilon
  + O(\|\epsilon\|^3).
\]
Squaring:
\[
  (f_\theta(x_i) - y_i)^2
  = (J_i^\top \epsilon)^2
  + 2(J_i^\top \epsilon) \cdot \tfrac{1}{2}\epsilon^\top \nabla^2_\theta f_{\theta^*}(x_i)\,\epsilon
  + O(\|\epsilon\|^4).
\]

For the cross term:
$\E\!\bigl[(J_i^\top \epsilon) \cdot \epsilon^\top \nabla^2 f_{\theta^*}(x_i)\,\epsilon\bigr]
= \sum_{a,b,c} (J_i)_a\,(\nabla^2 f)_{bc}\,\E[\epsilon_a\,\epsilon_b\,\epsilon_c] = 0$,
since all odd moments of the Gaussian distribution vanish.

For the leading term:
$\E[(J_i^\top \epsilon)^2] = \sigma^2\,\|J_i\|^2 = \sigma^2\, h_i$.

Therefore
$\E[\ell_i(\theta)] = \frac{\sigma^2\, h_i}{2} + O(\sigma^4)$,
and
\[
  \E_Q[L]
  = \frac{\sigma^2}{2n}\sum_{i=1}^n h_i + O(\sigma^4)
  = \frac{\sigma^2\,\cC(\theta^*)^2}{2} + O(\sigma^4).
\]

\textbf{Step~3 (KL divergence).}\;
The KL divergence between two Gaussians is given by the standard formula:
\begin{align*}
  \KL\!\bigl(\mathcal{N}(\theta^*, \sigma^2 I) \,\|\,
  \mathcal{N}(0, \sigma_0^2 I)\bigr)
  &= \frac{d}{2}\!\left(\frac{\sigma^2}{\sigma_0^2}
  - 1 - \ln\frac{\sigma^2}{\sigma_0^2}\right)
  + \frac{\|\theta^*\|^2}{2\sigma_0^2}.
\end{align*}

\textbf{Step~4 (Assembly).}\;
Substituting Steps~2 and~3 into the PAC-Bayes inequality of Step~1 yields
\eqref{eq:pb}.
\end{proof}

\subsection{Proof of Corollary~\ref{cor:pb} (PAC-Bayes for SGD-selected solutions)}

\begin{proof}
By \Cref{thm:main}~(II), $\cC(\theta^*)^2 \leq 2/\eta$ for any
$\theta^* \in \cM_\eta$.  Setting $\sigma^2 = \alpha\eta$ in \Cref{thm:pb}:
\[
  \E_Q[L]
  = \frac{\alpha\eta}{2}\,\cC(\theta^*)^2 + O(\alpha^2 \eta^2)
  \leq \frac{\alpha\eta}{2} \cdot \frac{2}{\eta} + O(\eta^2)
  = \alpha + O(\eta^2).
\]
The KL term becomes
$\frac{d}{2}\!\left(\frac{\alpha\eta}{\sigma_0^2} - 1
- \ln\frac{\alpha\eta}{\sigma_0^2}\right) + \frac{\|\theta^*\|^2}{2\sigma_0^2}$.
When $\alpha\eta \ll \sigma_0^2$, we have
$\frac{\alpha\eta}{\sigma_0^2} \approx 0$ and
$-\ln\frac{\alpha\eta}{\sigma_0^2} = \ln\frac{\sigma_0^2}{\alpha\eta} \gg 1$,
so the KL term is dominated by $\frac{d}{2}\ln\frac{\sigma_0^2}{\alpha\eta}$.
The full bound then scales as
$\alpha + O(\eta^2) + O\!\left(\sqrt{\frac{d\,\ln(1/\eta)}{n}}\right)$.
\end{proof}

%% =================================================================
\section{Reparameterization Invariance}\label{sec:reparam}
%% =================================================================

\begin{remark}[Reparameterization invariance of $\Sigma/H$]\label{rem:reparam}
  The ratio $\Sigma/H$ (and the interpolation-regime analogue $\cC^2/h_{\max}$)
  is invariant under smooth reparameterization.  This distinguishes our complexity
  measure from flatness (small $H$) or small noise ($\Sigma$) individually, both of
  which can be altered by reparameterization without changing the function
  \cite{Dinh2017}.

  \textbf{Proof.}\;
  Let $\theta = g(\varphi)$ be a smooth bijection with $\varphi^* = g^{-1}(\theta^*)$.
  In the new parameterization:
  \[
    \tilde{H} = L''(g(\varphi^*)) \cdot [g'(\varphi^*)]^2 = H \cdot (g')^2,
  \]
  \[
    \tilde{\Sigma}
    = \E_i\!\bigl[(\ell_i'(g(\varphi^*)))^2\bigr] \cdot [g'(\varphi^*)]^2
    = \Sigma \cdot (g')^2.
  \]
  Therefore $\tilde{\Sigma}/\tilde{H} = \Sigma \cdot (g')^2 / (H \cdot (g')^2)
  = \Sigma / H$.

  In the multivariate case, $\tr(H^{-1}\Sigma)$ is similarly invariant: the Jacobian
  factors from the change of variables cancel in the product $H^{-1}\Sigma$.

  \textbf{Implication.}\;
  The generalization gap $\approx \Sigma/(H(n-1))$ in the non-interpolation regime is
  an intrinsic geometric quantity.  In the interpolation regime, the per-sample
  curvature condition $h_{\max} < 2/\eta$ is coordinate-dependent (since $h_i$ involves
  the parameter-space norm of $J_i$), but the \emph{stability criterion} (whether
  the Lyapunov exponent is positive or negative) is invariant: a reparameterization
  that changes $h_i$ also changes the effective learning rate in the new coordinates,
  preserving the sign of $\lambda_k$.
\end{remark}

%% =================================================================
\section{Discussion}\label{sec:discussion}
%% =================================================================

\subsection{Connection to the non-interpolation regime}\label{ssec:non-interp}

In the non-interpolation regime ($L(\theta^*) > 0$), the gradient noise
$\Sigma(\theta^*) > 0$ enables an SDE approximation whose Fokker--Planck
stationary distribution $p \propto (1/\Sigma)\exp(-2V_{\mathrm{eff}}/\tau)$
biases SGD toward minima with low $\Sigma$.  At interpolation
($L(\theta^*) = 0$), $\Sigma = 0$ and the SDE degenerates; the
regularization mechanism shifts from first-order noise to the second-order
structure $H_i = J_i J_i^\top$ analyzed in this paper.  Both mechanisms
reflect per-sample sensitivity of $f_\theta$, but through different
mathematical tools (continuous SDE vs.\ discrete random matrix products).

\subsection{Testable predictions}

\begin{corollary}[Testable predictions from the theory]\label{cor:testable}
  Under the assumptions of \Cref{thm:main}:
  \begin{enumerate}[label=(G\arabic*)]
    \item \label{pred:hmax}
      \textbf{$h_{\max}$ predicts generalization.}\;
      After training to interpolation, measure
      $h_{\max} = \max_i \|\nabla_\theta f_{\theta^*}(x_i)\|^2$.
      Models trained with different $\eta$ should exhibit a positive correlation
      between $h_{\max}$ and test error, stronger than the correlation with
      $\tr(H)$ or other flatness measures.

    \item \label{pred:phase}
      \textbf{Critical phase transition.}\;
      For a fixed model, training loss should transition sharply from zero to nonzero
      at $\eta \approx 2/h_{\max}$.

    \item \label{pred:eta-mono}
      \textbf{Monotone effect of $\eta$.}\;
      In the interpolation regime, increasing $\eta$ should monotonically decrease
      $\cC = \|\bJ\|_F/\sqrt{n}$ (up to the point where SGD fails to converge).
  \end{enumerate}
\end{corollary}

\begin{conjecture}[Beyond the orthogonality assumption]\label{conj:beyond}
  \begin{enumerate}[label=(G\arabic*)]
    \setcounter{enumi}{2}
    \item \label{pred:noise-struct}
      \textbf{Noise structure matters more than magnitude.}\;
      Replacing SGD noise with isotropic Gaussian noise of the same magnitude should
      degrade generalization performance, because the structured per-sample filtering
      mechanism (which depends on the alignment of $J_i$ directions) is destroyed.

    \item \label{pred:F4}
      \textbf{F4 observability.}\;
      In small-data, large-model settings, overfitting solutions may have
      \emph{lower} $\Sigma$ than generalizing solutions, because memorization can
      lead to aligned per-sample gradients.
  \end{enumerate}
\end{conjecture}

\subsection{Limitations}

We discuss the main limitations of the present analysis.

\paragraph{Jacobian orthogonality (\Cref{A:orth}).}
The complete decoupling of coordinate dynamics relies on the orthogonality
$J_i^\top J_j = 0$.  While this is approximately satisfied in the NTK regime
(after whitening), real neural networks need not operate in this regime.
Removing \Cref{A:orth} requires the theory of products of non-commuting random
matrices (Furstenberg--Kesten theory), which yields qualitatively similar conclusions
but with less explicit constants.

\paragraph{Linearization gap.}
\Cref{lem:lin} establishes that the SGD dynamics near $\theta^*$ are
well-approximated by the linear system $\delta_{t+1} = A_{z_t}\,\delta_t$ only when
$\|\delta_t\|$ is small.  The stability conclusions hold in a neighborhood of
$\theta^*$; the size of this neighborhood depends on the ratio of the linearization
error $C\eta\|\delta\|^2$ to the linear term.  A complete quantitative analysis of
the basin of attraction requires nonlinear stability theory.

\paragraph{Behavior in $\cS^\perp$.}
The linearized SGD has no force in the orthogonal complement $\cS^\perp$: perturbations
in this subspace remain constant.  The full nonlinear dynamics may drift in $\cS^\perp$
due to higher-order effects, but characterizing this drift requires going beyond the
linearization.

\paragraph{Global trajectory.}
Both \Cref{thm:instability} and \Cref{thm:stability} are local stability results.
The global trajectory of SGD from random initialization to an element of $\cM_\eta$
involves traversing a complex loss landscape, and its analysis remains open.  Our
results establish \emph{which} solutions SGD can stably reside at, not \emph{how} it
gets there.

\subsection{Open problems}

\begin{enumerate}
  \item \textbf{Removing the orthogonality assumption.}\;
    Establish the instability threshold $\eta h_{\max} > 2$ for general (non-orthogonal)
    Jacobian configurations using Furstenberg's theory, with quantitative bounds on the
    Lyapunov exponent.

  \item \textbf{Global convergence.}\;
    Characterize the full trajectory of SGD from random initialization to $\cM_\eta$,
    combining the transient (non-interpolation) phase with the local stability analysis.

  \item \textbf{Discrete--continuous unification.}\;
    Develop a unified mathematical framework that handles the $\Sigma \to 0$ transition
    between the non-interpolation and interpolation regimes.

  \item \textbf{Architecture-dependent analysis.}\;
    Determine for which architectures and data distributions the condition
    ``low $h_{\max}$ implies good generalization'' (the analogue of $\neg F4$ in the
    notation of Section~\ref{ssec:non-interp}) is satisfied.

  \item \textbf{Mini-batch and momentum.}\;
    Extend the per-sample Lyapunov analysis to mini-batch SGD and SGD with momentum,
    determining how the critical learning rate $\eta_c$ is modified.
\end{enumerate}

%% =================================================================
\section{Related Work}\label{sec:related}
%% =================================================================

\paragraph{Implicit regularization of SGD.}
Smith and Le~\cite{SmithLe2018} observed that larger learning rates and smaller batch
sizes improve generalization, and proposed the SDE approximation
$d\theta = -\nabla L\,dt + \sqrt{(\eta/B)\Sigma}\,dW_t$ to explain this
phenomenon.  Chaudhari and Soatto~\cite{ChaudhariSoatto2018} analyzed the
Fokker--Planck equation of this SDE and showed that the stationary distribution
$p \propto (1/\Sigma)\exp(-2V_{\mathrm{eff}}/\tau)$ biases SGD toward minima with low
noise covariance.

\paragraph{Edge of stability and sharpness reduction.}
Wu, Wang, and Su~\cite{WuWangSu2022} studied the phenomenon of sharpness reduction
during training, where the largest eigenvalue of the Hessian tends to decrease toward
$2/\eta$.  Wu and Su~\cite{WuSu2023} proved that SGD cannot stably converge to
solutions where $\tr(H) > 2/\eta$.  Our work refines this by showing that the relevant
condition is per-sample: $h_{\max} > 2/\eta$ (not the aggregate $\tr(H) > 2/\eta$),
and by closing the loop from stability to generalization.  The distinction matters
because $\tr(H) = \sum_i h_i / n \cdot n = \sum_i h_i$ (for appropriately defined
$H$), and $\tr(H) \leq 2/\eta$ is compatible with having one very large $h_k$
compensated by many small ones, whereas our condition is not.

\paragraph{Reparameterization invariance.}
Dinh et al.~\cite{Dinh2017} demonstrated that Hessian-based sharpness measures are
not invariant under reparameterization and constructed explicit examples where
``sharp'' minima (in terms of Hessian eigenvalues) generalize well.  Our
\Cref{rem:reparam} shows that the ratio $\Sigma/H$ (and its interpolation-regime
analogue) is reparameterization-invariant, providing a more robust complexity measure.

\paragraph{Li et al.\ (2021).}
Li et al.~\cite{Li2021} studied the implicit bias of SGD toward flat minima in the
context of label noise SGD, providing theoretical evidence that SGD noise structure
matters for generalization.

\paragraph{Chang and Khanna (2025).}
Chang and Khanna~\cite{ChangKhanna2025} studied the relationship between SGD dynamics
and generalization in overparameterized models, providing complementary perspectives
on the role of stochastic noise.

%% =================================================================
%% REFERENCES
%% =================================================================
\begin{thebibliography}{10}

\bibitem{ChaudhariSoatto2018}
P.~Chaudhari and S.~Soatto.
\newblock Stochastic gradient descent performs variational inference, converges to
  limit cycles for deep networks.
\newblock In \emph{Proceedings of ICLR}, 2018.

\bibitem{Dinh2017}
L.~Dinh, R.~Pascanu, S.~Bengio, and Y.~Bengio.
\newblock Sharp minima can generalize for deep nets.
\newblock In \emph{Proceedings of ICML}, pages 1019--1028, 2017.

\bibitem{LiTaiE2017}
Q.~Li, C.~Tai, and W.~E.
\newblock Stochastic modified equations and adaptive stochastic gradient algorithms.
\newblock In \emph{Proceedings of ICML}, pages 2101--2110, 2017.

\bibitem{Li2021}
Z.~Li, T.~Wang, and S.~Arora.
\newblock What happens after {SGD} reaches zero loss?---{A} mathematical framework.
\newblock In \emph{Proceedings of ICLR}, 2021.

\bibitem{SmithLe2018}
S.~L. Smith and Q.~V. Le.
\newblock A {B}ayesian perspective on generalization and stochastic gradient descent.
\newblock In \emph{Proceedings of ICLR}, 2018.

\bibitem{WuSu2023}
L.~Wu and W.~J. Su.
\newblock The implicit regularization of dynamical stability in stochastic gradient
  descent.
\newblock In \emph{Proceedings of ICML}, 2023.

\bibitem{WuWangSu2022}
L.~Wu, X.~Wang, and W.~J. Su.
\newblock Does the sharpness of SGD iterates control generalization?
\newblock \emph{arXiv preprint arXiv:2207.06680}, 2022.

\bibitem{ChangKhanna2025}
H.~Chang and R.~Khanna.
\newblock Implicit regularization of stochastic gradient descent in overparameterized
  models.
\newblock \emph{arXiv preprint}, 2025.

\bibitem{McAllester1999}
D.~A. McAllester.
\newblock {PAC-Bayesian} model averaging.
\newblock In \emph{Proceedings of COLT}, pages 164--170, 1999.

\bibitem{Catoni2007}
O.~Catoni.
\newblock \emph{{PAC-Bayesian} Supervised Classification: The Thermodynamics of
  Statistical Learning}.
\newblock IMS Lecture Notes, vol.~56. Institute of Mathematical Statistics, 2007.

\bibitem{MohriRostamizadehTalwalkar2018}
M.~Mohri, A.~Rostamizadeh, and A.~Talwalkar.
\newblock \emph{Foundations of Machine Learning}.
\newblock MIT Press, 2nd edition, 2018.

\bibitem{furstenberg1960products}
H.~Furstenberg and H.~Kesten.
\newblock Products of random matrices.
\newblock \emph{Annals of Mathematical Statistics}, 31(2):457--469, 1960.

\end{thebibliography}

\end{document}
